\chapter{Einstein Gravity Plus a Conserved Extra Tensor}
\label{app:field-eq}

This appendix writes the infrared field equation of
Theorems~\ref{thm:mT}--\ref{thm:conservation} next to
Einstein's equation. It introduces no new hypothesis.

\section{Einstein}
\begin{equation}
  \label{eq:appF-Ein}
  G_{\mu\nu}+\Lambda g_{\mu\nu}
  =8\pi G\,T_{\mu\nu}^{\mathrm{matter}}.
\end{equation}
Bianchi identities force
\(\nabla^\mu T_{\mu\nu}^{\mathrm{matter}}=0\) on shell.

\section{NAMCG infrared equation}
After the composite metric, \(m_T=0\), the heat-kernel map,
Lorentz locking, and conservation are granted,
\begin{equation}
  \label{eq:appF-NAMCG}
  G_{\mu\nu}+\Lambda_{\mathrm{eff}}g_{\mu\nu}
  =8\pi\Geff
  \bigl(
    T_{\mu\nu}^{\mathrm{matter}}
   +T_{\mu\nu}^{\mathrm{tensor}}
   +T_{\mu\nu}^{\mathrm{ent}}
  \bigr).
\end{equation}
\(G_{\mu\nu}\) is the Einstein tensor of the composite metric
\(g_{\mu\nu}=\eta_{AB}e^A{}_\mu e^B{}_\nu\).
\(\Lambda_{\mathrm{eff}}\) is the coefficient of
Theorem~\ref{thm:Leff}; \(G_{\mathrm{eff}}\) is to be matched by the
conditional sector calculation of Problem~F.
The extra source is
\begin{equation}
  \label{eq:appF-Tent}
  T_{\mu\nu}^{\mathrm{ent}}
  =-\frac{2}{\sqrt{-g}}
   \frac{\delta U_{\mathrm{ent}}[I(A:B)]}{\delta g^{\mu\nu}}
  =T_{\mu\nu}^{\mathrm{ent,rev}}
   +T_{\mu\nu}^{\mathrm{ent,irrev}}.
\end{equation}
Chapter~11 requires
\begin{equation}
  \label{eq:appF-cons}
  \nabla^\mu
  \bigl(
    T_{\mu\nu}^{\mathrm{matter}}
   +T_{\mu\nu}^{\mathrm{tensor}}
   +T_{\mu\nu}^{\mathrm{ent}}
  \bigr)
  =0.
\end{equation}
In the infrared, \(T_{\mu\nu}^{\mathrm{tensor}}\) is absorbed
into the Einstein tensor of \(g_{\mu\nu}\). What can appear
as extra gravity in static disks is
\(T_{\mu\nu}^{\mathrm{ent,rev}}\).

\section{Weak field}
Newtonian gauge, static source:
\begin{equation}
  \label{eq:appF-Poisson}
  \nabla^2\Phi
  =4\pi\Geff
   \bigl(\rho_{\mathrm{b}}+\rho_{\mathrm{ent}}\bigr),
  \qquad
  \rho_{\mathrm{ent}}
  =T_{00}^{\mathrm{ent,rev}}.
\end{equation}
Assumption~\ref{ass:track} slaved \(\rho_{\mathrm{ent}}\) to
\(\rho_{\mathrm{b}}\). The interpolator \(\nu_{\sqrt{}}\) of
Appendix~\ref{app:empirical} is a coefficient choice for that
assumption, not a consequence of~\eqref{eq:appF-NAMCG}.

\section{What this does and does not change}
High acceleration:
\(T_{\mu\nu}^{\mathrm{ent,rev}}\to 0\), so~\eqref{eq:appF-NAMCG}
returns to~\eqref{eq:appF-Ein} with
\((\Geff,\Lambda_{\mathrm{eff}})\) in place of \((G,\Lambda)\).
Low acceleration: \(\rho_{\mathrm{ent}}\) may dominate
\(\rho_{\mathrm{b}}\) without a particle halo, if and only if
Assumption~\ref{ass:track} holds.
Clusters, the CMB, and colliding systems test that assumption,
not \(G_{\mu\nu}\).
Gravitational waves test \(m_T=0\), not~\eqref{eq:appF-Tent}.
The irreversible drag of Chapter~14 is a phenomenological addition,
not a static extra potential and not derived from
\eqref{eq:appF-Tent}.

\section{One-line summary}
Einstein equates curvature to ordinary stress. NAMCG equates
the same curvature to ordinary stress plus a divergenceless
entanglement stress of baryonic regions.